% sections/00-introduction.tex
\section{Introduction}
\label{sec: introduction}

\subsection{The geometry of the configuration group}

The configuration space of an ideal incompressible fluid on a closed Riemannian
manifold $M$ is the group $\SDiff(M)$ of volume-preserving diffeomorphisms,
carrying the right-invariant $L^2$ metric. Arnold showed that the geodesics of
this metric are the solutions of the Euler equations~\cite{arnold-1966-sur-la-geo},
and Ebin and Marsden placed the resulting infinite-dimensional geometry on a
rigorous analytic footing~\cite{ebin-1970-gro-dif-mot}. On the Lie algebra
$\Lsig$ of smooth divergence-free vector fields the Levi-Civita connection of the
$L^2$ metric is the projection of the flat ambient derivative,
\begin{equation}\label{eq: connection}
  \conn_u v = \Leray\bigl[(u\cdot\nabla)v\bigr],
\end{equation}
with $\Leray$ the Leray projection onto divergence-free fields. Arnold computed
the sectional curvatures of~\eqref{eq: connection} and found them predominantly
negative~\cite{arnold-1966-sur-la-geo, arnold-1998-top-met-hyd}. Negative
sectional curvature forces the exponential divergence of nearby geodesics through
the Jacobi equation, and so measures the Lagrangian instability of the flow; the
equivalence of this geometric instability with the Eulerian instability of the
velocity field was established by Rouchon, Misio\l{}ek and
Preston~\cite{rouchon-1992-jac-cur, misiolek-1993-sta-flo-ide,
preston-2004-ide-flu-ins}. The sign of these curvatures is carried by the
pressure: for a pair of divergence-free fields the sectional value is, modulo
terms that vanish on Fourier modes, the negative $L^2$ energy of an associated
pressure gradient~\cite{arnold-1998-top-met-hyd}, so the incompressibility
constraint is the source of the curvature.

\subsection{The regularity problem}

For the viscous problem on $\T^3$, Leray and Hopf constructed global weak
solutions of finite energy~\cite{leray-1934-sur-le-mou, hopf-1951-ube-die-anf},
and whether every smooth datum gives a globally smooth solution is the subject of
the Clay millennium problem~\cite{fefferman-2006-exi-smo}. Smoothness follows
from a range of scaling-critical conditions on the velocity, the
Prodi--Serrin--Ladyzhenskaya class~\cite{prodi-1959-un-teo-di,
serrin-1962-int-reg-wea, ladyzhenskaya-1967-uni-smo-gen}, with the endpoint
$L^\infty_t L^3_x$ secured by Escauriaza, Seregin and
\v{S}ver\'ak~\cite{escauriaza-2003-l3i-nav-equ}; the vorticity counterpart for
the inviscid problem is the criterion of Beale, Kato and
Majda~\cite{beale-1984-rem-bre-smo}. Where global smoothness is open, the
partial-regularity theory of Scheffer and of Caffarelli, Kohn and Nirenberg
bounds the singular set, which has parabolic Hausdorff dimension at most
one~\cite{scheffer-1976-par-reg-sol, caffarelli-1982-par-reg-sui}, with later
proofs given by Lin and by Vasseur~\cite{lin-1998-new-pro-caf,
vasseur-2007-new-pro-par}.

\subsection{Vortex stretching and the pressure Hessian}

The mechanism behind these results is vortex stretching, the growth of vorticity
driven by the scalar $\alpha = \xi^{\top} S\,\xi$, with $\xi$ the vorticity
direction and $S$ the strain. Constantin and Fefferman showed that a solution
whose vorticity direction stays coherent on the set of high vorticity remains
smooth~\cite{constantin-1993-dir-vor-pro}, a geometric depletion of the
nonlinearity developed further in the geometric statistics of turbulence and the
constraints on singular Euler flows~\cite{constantin-1994-geo-sta-tur,
constantin-1996-geo-con-pot} and in the regularising effect of the vorticity
direction for the viscous problem~\cite{beiraodaveiga-2002-reg-eff-vor}. The
strain that drives $\alpha$ evolves through the pressure Hessian $\nabla^2 p$, a
nonlocal transform of the velocity gradient whose intractability is the central
difficulty of the local models of the velocity-gradient tensor introduced by
Vieillefosse and Cantwell, and of the analysis of vortex stretching by Ohkitani
and Kishiba~\cite{vieillefosse-1982-loc-int, cantwell-1992-exa-sol,
ohkitani-1995-non-nat}. The geometry of the configuration group and the analysis
of vortex stretching have thus developed alongside one another, the first through
the curvature of~\eqref{eq: connection}, the second through the pressure Hessian.
The purpose of this paper is to show that these are one object.

\subsection{Results}

Throughout, $u$ is smooth and divergence-free on the flat torus $\T^3$, with
strain $S_{ij} = \tfrac12(\partial_i u_j + \partial_j u_i)$, vorticity
$\omega = \curl u$, vorticity direction $\xi = \omega/\abs{\omega}$ on
$\{\omega \neq 0\}$, and pressure $p$ determined by
\begin{equation}\label{eq: pressure-poisson}
  -\Delta p = \partial_i u_j\,\partial_j u_i = \abs{S}^2 - \tfrac12\abs{\omega}^2
            =: q.
\end{equation}
We write $\nabla^2 f$ for the Hessian of a function $f$, and $\Riesz$ for the
matrix Riesz transform with symbol $\eta_i\eta_j/\abs{\eta}^2$, so that
$\nabla^2 p = -\Riesz q$.

The contribution is a reading of the Navier--Stokes singular set through the
curvature of the configuration group. At a singular point of a suitable weak
solution the partial-regularity theory leaves a dichotomy
(Theorem~\ref{thm: dichotomy}): the curvature concentrates, or the flow comes into
strain--enstrophy balance. Away from balance the curvature controls regularity
coercively (Proposition~\ref{prop: coercive}), and in the scale-critical class the
balanced alternative is excluded, so every scale-critical singular point, the
Type-I points among them, is curvature-concentrating (Section~\ref{sec: blowup}).
Beyond that class the balanced regime is supercritical, where the energy identity is
blind and the curvature is the structure a regularity proof must read
(Section~\ref{sec: barrier}). This reading rests on a geometric identity, that the
curvature of $\SDiff(\T^3)$ is the pressure Hessian, which we establish first.

That the sectional curvature of $\SDiff(\T^3)$ is carried by the pressure is due to
Arnold~\cite{arnold-1966-sur-la-geo, arnold-1998-top-met-hyd}, with the analytic
foundation of Ebin and Marsden~\cite{ebin-1970-gro-dif-mot} and the curvature
computations of Rouchon, Misio\l{}ek and Preston~\cite{rouchon-1992-jac-cur,
misiolek-1993-sta-flo-ide, preston-2004-ide-flu-ins}. The identity established
here is at the level of the curvature operator: $\curv(u,\cdot)u$ is the
Leray-projected pressure Hessian $-\Leray[(\nabla^2 p)\,\cdot]$, the same operator
$\nabla^2 p = -\Riesz q$ that drives the velocity-gradient dynamics of
Vieillefosse, Cantwell, Ohkitani, Meneveau and Johnson~\cite{vieillefosse-1982-loc-int,
cantwell-1992-exa-sol, ohkitani-1995-non-nat, meneveau-2011-lag-dyn-mod,
johnson-2021-rol-vor-str}. This identifies the configuration-space curvature with
the pressure Hessian of that dynamics, measures it by the strain--enstrophy
imbalance, and turns it on the singular set through a transport law and a blow-up
limit. The three theorems follow, one to a section, and the dichotomy is read
through the dynamics. The geometry of Sections~\ref{sec: connection-curvature}--\ref{sec: imbalance},
the velocity-gradient dynamics it identifies, and the partial-regularity and
divergence-free-drift Liouville theory that the singular-set reading draws on are
established. The contribution is the operator identity that binds the geometry to
the dynamics, and the reading of the singular set it gives.

Theorem~\ref{thm: main} (Section~\ref{sec: connection-curvature}) identifies the
curvature with the pressure Hessian. The second fundamental form of
$\SDiff(\T^3)$ in the flat group $\Diffop(\T^3)$ is the pressure gradient, the
flat ambient part of the curvature vanishes, and
$\curv(u,v)u = -\Leray[(\nabla^2 p)\,v] - \Leray[(\nabla^2 q_v)\,u]$, with $q_v$
the mixed pressure $\Delta q_v = \partial_i v_j\,\partial_j u_i$; the Gauss
equation reproduces Arnold's sectional curvature formula. This is the Gauss
equation for the constraint of incompressibility: the ambient group has no
curvature, and the whole curvature of $\SDiff(\T^3)$ is carried by the second
fundamental form, which is the pressure.

Theorem~\ref{thm: imbalance} (Section~\ref{sec: imbalance}) measures it. The
pointwise trace of the curvature is the strain--enstrophy imbalance,
$\operatorname{tr}\nabla^2 p = \Delta p = -q$, and its $L^2$ size is that
imbalance, $\norm{\nabla^2 p}_{L^2(\T^3)} = \norm{q}_{L^2(\T^3)}$. The curvature
is large where strain and enstrophy are out of balance and vanishes on the
pressureless flows $\{q \equiv 0\}$, the single-field form of the same-shell
degeneracy that Arnold found for Fourier-mode pairs.

Theorem~\ref{thm: dichotomy} (Section~\ref{sec: dichotomy}) reads this against
regularity. Because $\abs{\nabla u}^2 = q + \abs{\omega}^2$ pointwise, at a
singular point of a suitable weak solution either the curvature concentrates,
$\limsup_{r\to0} r^{-1}\!\int_{Q_r}\abs{q} > 0$, or the scaled strain energy and
half-enstrophy converge to a common value: the flow comes into strain--enstrophy
balance, equivalently it becomes locally pressureless. Away from balance the
converse holds: the curvature controls the local energy coercively
(Proposition~\ref{prop: coercive}), so the two directions place the singular set
where the curvature concentrates or the pressure depletes. The
Constantin--Fefferman criterion controls the stretching scalar
$\alpha = \xi^{\top} S\,\xi$ under a coherence hypothesis on $\xi$ and reaches the
regime of coherent vorticity; the curvature, through its imbalance, governs the
complementary regime.

Two further results read the balanced alternative dynamically.
Section~\ref{sec: dichotomy} gives the imbalance a transport law,
$(\partial_t + u\cdot\nabla)q = -6\det(\nabla u) - 2\,S\!:\!\nabla^2 p
+ 2\nu\operatorname{tr}(\nabla u\,\Delta\nabla u)$, in which each term carries no
net source, so balance demands their pointwise cancellation against a zero-sum
budget. Section~\ref{sec: blowup} makes this rigid: a balanced singular point would
produce a bounded ancient solution with no pressure, $q\equiv0$ and
$\nabla p\equiv0$. Under a Type-I bound the limit inherits the self-similar decay
$\abs{U}\le M\abs{t}^{-1/2}$ and an elementary maximum principle forces it to
vanish; under the larger scale-critical $L^3$ bound the divergence-free drift lies
in the critical class and the Liouville theorem of Seregin, Silvestre, \v{S}ver\'ak
and Zlato\v{s}~\cite{seregin-2012-div-dri} forces the same. Every scale-critical
singular point is therefore curvature-concentrating. Beyond the critical class the
balanced regime is supercritical, and the averaged-equation barrier of
Tao~\cite{tao-2016-fin-tim-blo} shows it is closed to the energy identity and the
scaling alone; the curvature is the structure of the Euler nonlinearity that
resolves it, and the reduced pressureless Liouville is the gateway it leaves open.

Section~\ref{sec: proof} proves Theorem~\ref{thm: main}; the remaining sections
are self-contained. The identities are confirmed by an independent symbolic and
spectral computation, described in Appendix~\ref{app: verification}.
